\documentclass[a4paper, 12pt]{extarticle}

% Поддержка языков
\usepackage[english, russian]{babel} 

% Настройка кодировок
\usepackage[T2A]{fontenc}
\usepackage[utf8]{inputenc}

% Настройка шрифтов
\usepackage{fontspec}
\setmainfont[Ligatures=TeX]{Times New Roman} % Шрифт для основного текста документа
\setsansfont[Ligatures=TeX]{Arial}
\setmonofont{Consolas} % Шрифт для кода

% Настройка отступов от краев страницы
\usepackage[left=3cm, right=1.5cm, top=2cm, bottom=2cm]{geometry}

\usepackage{titleps} % Колонтитулы
\usepackage{subfig} % Для подписей к рисункам и таблицам
\usepackage{graphicx} % для вставки картинок
\graphicspath{{./img/}} % Путь до папки с изображениями
% Пакет для отрисовки графиков
\usepackage{tikz}
\usetikzlibrary{arrows,positioning,shadows}
\usepackage{stmaryrd} % Стрелки в формулах
\usepackage{indentfirst} % Красная строка после заголовка
\usepackage{hhline} % Улучшенные горизонтальные линии в таблицах
\usepackage{multirow} % Ячейки в несколько строчек в таблицах
\usepackage{longtable} % Многостраничные таблицы
\usepackage{paralist,array} % Список внутри таблицы
\usepackage[normalem]{ulem}  % Зачеркнутый текст
\usepackage{upgreek, tipa} % Красивые греческие буквы
\usepackage{amsmath, amsfonts, amssymb, amsthm, mathtools} % ams пакеты для математики, табуляции
\usepackage{nicematrix} % Особые матрицы pNiceArray

\linespread{1.5} % Межстрочный интервал
\setlength{\parindent}{1.25cm} % Табуляция
\setlength{\parskip}{0cm}

% Пакет для красивого выделения кода
\usepackage{minted}
\setminted{fontsize=\footnotesize}

% Добавляем гипертекстовое оглавление в PDF
\usepackage[
bookmarks=true, colorlinks=true, unicode=true,
urlcolor=black,linkcolor=black, anchorcolor=black,
citecolor=black, menucolor=black, filecolor=black,
]{hyperref}

% Убрать переносы слов
\tolerance=1
\emergencystretch=\maxdimen
\hyphenpenalty=10000
\hbadness=10000

\newpagestyle{main}{
	% Верхний колонтитул
	\setheadrule{0cm} % Размер линии отделяющей колонтитул от страницы
	\sethead{}{}{} % Содержание {слева}{по центру}{справа}
	% Нижний колонтитул
	\setfootrule{0cm} % Размер линии отделяющей колонтитул от страницы
	\setfoot{}{}{\thepage} % Содержание {слева}{по центру}{справа}
}
\pagestyle{main}

% НОВЫЕ КОМАНДЫ
\newcommand{\deriv}[2]{\frac{\partial #1}{\partial #2}}
\newcommand{\n}{\par}
\newcommand{\percent}{\mathbin{\%}}

% Заменяем Рис. на Рисунок
\addto\captionsrussian{\renewcommand{\figurename}{Рисунок}}

% Изменение формата подписей
% Стиль номера таблицы/рисунка #1-Таблица/Рис. #2-номер
\DeclareCaptionLabelFormat{custom}{#1 #2}
% Стиль разделителя номера таблицы/рисунка и названия таблицы/рисунка
\DeclareCaptionLabelSeparator{custom}{ --- }
% Стиль формата #1-номер таблицы/рисунка #2-разделитель #3-название
\DeclareCaptionFormat{custom}{#1#2#3}


\captionsetup
{
	format=custom,%
	labelsep=custom,
	labelformat=custom
}

% ПЕРЕГРУЗКА УЖЕ СУЩЕСТВУЮЩИХ КОМАНД
\renewcommand{\epsilon}{\varepsilon} % Заменить знак эпсилон
\renewcommand{\phi}{\varphi}
\renewcommand{\kappa}{\varkappa}
\renewcommand{\lambda}{\uplambda}


\usepackage{float}
 

\title{\textbf{Понятие о межсетевом экранировании. Компоненты межсетевого экрана. \\ Политика сетевой безопасности. Критерии фильтрации пакетов. \\ Основные схемы сетевой защиты на базе межсетевых экранов. \\ Требования к межсетевым экранам. \\ Списки доступа и их настройка.}}	
\date{}
	
\begin{document}
	\maketitle
	\thispagestyle{empty}
	\newpage
	
	\section*{Понятие о межсетевом экранировании}
	
	\subsection*{Краткое определение межсетевого экрана}
	\begin{itemize}
		\item 	Межсетевой экран (Firewall, брендмауэр) - это программное или программно-аппаратное средство, осуществляющее контроль и регламентацию передачи сетевых потоков в соответствии с заданной политикой безопасности.
	\end{itemize}

	\subsection*{Полное определение межсетевого экрана}
	\begin{itemize}
		\item Межсетевой экран - это локальное (однокомпонентное) или функционально распределенное программное (программно-аппаратное) средство (комплекс), реализующее контроль за информацией, поступающей в АС и/или выходящей из АС.
		\item МЭ обеспечивает защиту АС посредством фильтрации информации, т.е. ее анализа по совокупности критериев и принятия решения о ее распространении в (из) АС на основе заданных правил, проводя таким образом разграничение доступа субъектов из одной АС к объектам другой АС.
		\item Каждое правило запрещает или разрешает передачу информации определенного вида между субъектами и объектами. Как следствие субъекты из одной АС получают доступ только к разрешенным информационным объектам из другой АС.
		\item Интерпретация набора правил выполняется последовательностью фильтров, которые разрешают или запрещают передачу данных (пакетов) на следующий фильтр или уровень протокола (РД Гостехкомиссии).
	\end{itemize}
	
	\subsection*{Угрозы, на противодействие которым направлены МЭ}
	
	\begin{itemize}
		\item Использование небезопасных протоколов и служб;
		\item Инвентаризация сетевых ресурсов;
		\item Несанкционированное информационное взаимодействие;
		\item Сетевые атаки (DoS, flood).
	\end{itemize}
	
	\subsection*{Функции межсетевого экрана}
	
	\begin{itemize}
		\item Фильтрация сетевого трафика на различных уровнях модели ISO/OSI;
		\item Контроль состояния протоколов, инспекция и нормализация сетевого трафика;
		\item Трансляция сетевых адресов;
		\item Маршрутизация и участие в протоколах динамической маршрутизации;
		\item Кеширование передаваемых данных.
	\end{itemize}
	
	\newpage
	
	\section*{Компоненты межсетевого экрана}
	 
	\begin{itemize}
		\item Фильтрующие маршрутизаторы;
		\item Шлюзы сеансового уровня;
		\item Шлюзы уровня приложений.
	\end{itemize}
	
	\newpage
	
	\section*{Политика сетевой безопасности}
	Политика сетевой безопасности (ГОСТ Р ИСО/МЭК 27033-1-2011) - это совокупность положений, правил и практических приемов, устанавливающих подход организации к использованию ее сетевых ресурсов и определяющих, как должна обеспечиваться защита ее сетевой инфраструктуры и сервисов.
	
	\begin{itemize}
		\item Безопасная передача данных:
		
		\begin{enumerate}
			\item Виртуальные частные сети;
			\item Безопасность сетевой инфраструктуры;
			\item Безопасность беспроводных сетей.
		\end{enumerate}
		
		\item Защита от сетевых атак:
		
		\begin{enumerate}
			\item Межсетевое экранирование;
			\item Антивирусная защита;
			\item Системы обнаружения атак;
			\item Защита серверов и рабочих станций;
			\item L2-безопасность.
		\end{enumerate}
		
		\item Идентификация и контроль пользователе:
		
		\begin{enumerate}
			\item Системы аутентификации;
			\item Системы контроля доступа;
			\item Системы управления учетными записями;
		\end{enumerate}
		
		\item Мониторинг и управление системой безопасности.
	\end{itemize}
	
	\section*{Критерии фильтрации пакетов}
	
	\subsection*{Фильтрация трафика}
	
	\begin{itemize}
		\item IP-адреса - фильтрация по адресу отправителя и получателя;
		\item Сетевые опции - параметры внутри заголовка пакета;
		\item Протокол транспортного уровня;
		\item Флаги и опции транспортного уровня;
		\item Порты - фильтрация по номерам портов отправителя и получателя.
	\end{itemize}
	
	\subsection*{Особенности фильтрации трафика}

	\begin{center}
		\begin{tabular}{|p{5cm}|p{7cm}|}
			\hline
			\textbf{Преимущества} & \textbf{Недостатки} \\
			\hline
			Быстро и просто & Почти не работает в текущее время \\
			\hline
		\end{tabular}
	\end{center}
	
    \newpage
    
    \section*{Основные схемы сетевой защиты на базе межсетевых экранов}
    
    \subsection*{МЭ - фильтрующий маршрутизатор}
    
    Межсетевой экран, основанный на фильтрации пакетов, является самым распространенным и наиболее простым в реализации. Он состоит из фильтрующего маршрутизатора, расположенного между защищаемой и внешней сетью. Фильтрующий маршрутизатор сконфигурирован для блокирования или фильтрации входящих и исходящих пакетов на основе анализа их адресов и портов. 
    
    Компьютеры, находящиеся в защищаемой сети, имеют прямой доступ во внешнюю сеть, в то время как большая часть доступа к ним из внешней сети блокируется. фильтрующий маршрутизатор может реализовать любую из политик безопасности, описанных ранее.
    
    \begin{figure}[H]
    	\centering
    	\includegraphics[width=1\textwidth]{1}
    	\caption{Фильтрующий маршрутизатор}
    	\label{fig:mesh1}
    \end{figure}
    
    
    \subsection*{МЭ на основе двупортового шлюза}
    
    Межсетевой экран на базе двупортового шлюза включает двудомный хост-компьютер с двумя сетевыми интерфейсами. При передаче информации между этими интерфейсами и осуществляется основная фильтрация. Для обеспечения дополнительной защиты между прикладным шлюзом и внешней сетью обычно размещают фильтрующий маршрутизатор. В результате между прикладным шлюзом и маршрутизатором образуется внутренняя экранированная подсеть. Эту подсеть можно использовать для размещения доступных извне информационных серверов.
    
    В отличие от схемы межсетевого экрана с фильтрующим маршрутизатором прикладной шлюз полностью блокирует трафик между внешней и защищаемой сетью. Только полномочные серверы-посредники, располагаемые на прикладном шлюзе, могут предоставлять услуги и доступ пользователям.
    
    Данный вариант межсетевого экрана реализует политику безопасности, основанную на принципе «запрещено все, что не разрешено в явной форме», при этом пользователю недоступны все службы, кроме тех, для которых определены соответствующие полномочия. Такой подход обеспечивает высокий уровень безопасности, поскольку маршруты к защищенной подсети известны только межсетевому экрану и скрыты от внешних систем.
    
    Следует отметить, что безопасность двудомного хост-компьютера, используемого в качестве прикладного шлюза, должна поддерживаться на высоком уровне. Любая брешь в его защите может серьезно ослабить безопасность защищаемой сети. Если шлюз окажется скомпрометированным, у злоумышленника появится возможность проникнуть в защищаемую сеть.
    
    Этот межсетевой экран может требовать от пользователей применения средств усиленной аутентификации для регистрации доступа, попыток зондирования и атак системы нарушителем.
    
    \begin{figure}[H]
    	\centering
    	\includegraphics[width=1\textwidth]{2}
    	\caption{Двупортовый шлюз}
    	\label{fig:mesh1}
    \end{figure}
    
    \subsection*{МЭ на основе экранирующего шлюза}
    
    В этой схеме первичная безопасность обеспечивается фильтрующим маршрутизатором. Пакетная фильтрация в фильтрующем маршрутизаторе может быть реализована одним из следующих способов:
    
    \begin{itemize}
    	\item позволять внутренним хост-компьютерам открывать соединения с хост-компьютерами во внешней сети для определенных сервисов;
   		\item запрещать все соединения от внутренних хост-компьютеров.
   	\end{itemize}
    
    Эти подходы можно комбинировать для различных сервисов, разрешая некоторым сервисам соединение непосредственно через пакетную фильтрацию, в то время как другим — только непрямое соединение через полномочные серверы-посредники. Все зависит от конкретной политики безопасности, принятой во внутренней сети. В частности, пакетная фильтрация на фильтрующем маршрутизаторе может быть организована таким образом, чтобы прикладной шлюз, используя свои полномочные серверы-посредники, обеспечивал для систем защищаемой сети доступ к различным сервисам.
    
    Межсетевой экран, выполненный по данной схеме, получается более гибким, но менее безопасным по сравнению с межсетевым экраном с прикладным шлюзом на базе двудомного компьютера. Это обусловлено тем, что в схеме межсетевого экрана с экранированным шлюзом существует потенциальная возможность передачи трафика в обход прикладного шлюза непосредственно к системам локальной сети.
    
    Основной недостаток заключается в том, что если атакующий нарушитель сумеет проникнуть в хост-компьютер, то перед ним окажутся незащищенные системы внутренней сети. Другой недостаток связан с возможной компрометацией маршрутизатора. Если маршрутизатор окажется скомпрометированным, внутренняя сеть станет доступна атакующему нарушителю.
    
    \begin{figure}[H]
    	\centering
    	\includegraphics[width=1\textwidth]{3}
    	\caption{Экранирующий шлюз}
    	\label{fig:mesh1}
    \end{figure}
    
	 \subsection*{МЭ - экранирующая подсеть}

	 Эта схема является развитием экранированного шлюза и использует два маршрутизатора для создания изолированной DMZ (демилитаризованной зоны) между внешней и внутренней сетью. В DMZ размещаются прикладной шлюз и публичные серверы.
	 
	 Внешний маршрутизатор защищает DMZ и внутреннюю сеть от внешней. Он пропускает трафик только к сервисам в DMZ и блокирует прямой доступ во внутреннюю сеть.
	 
	 Внутренний маршрутизатор защищает локальную сеть как от внешних атак, так и от скомпрометированных систем в DMZ. Он разрешает трафик только через доверенные узлы DMZ.
	 
	 Для проникновения во внутреннюю сеть злоумышленнику необходимо преодолеть два уровня защиты. Даже при взломе сервера в DMZ внутренняя сеть остается защищенной внутренним фильтром.
	 
	 Основной недостаток заключается в том, что пара фильтрующих маршрутизаторов нуждается в большом внимании для обеспечения необходимого уровня безопасности, поскольку из-за ошибок при их конфигурации могут возникнуть провалы в безопасности всей сети.
	 
	 \begin{figure}[H]
	 	\centering
	 	\includegraphics[width=1\textwidth]{4}
	 	\caption{Экранирующая подсеть}
	 	\label{fig:mesh1}
	 \end{figure}
	 
	 \newpage
	 
	\section*{Требования к межсетевым экранам}	 
	\begin{itemize}
			\item Функциональность - возможность работы на разных уровнях модели OSI, поддержка современных протоколов и сервисов;
			\item Пропускная способность - максимальный объем данных, который может быть передан за определенный промежуток времени;
			\item Поддержка технологии QoS - механизма управления сетевым трафиком, который позволяет выделять приоритеты для разных типов данных;
			\item Управляемость - наличие удобных интерфейсов конфигурации и возможность централизованного управления множеством устройств из единой консоли;
			\item Отказоустойчивость - способность системы сохранять работоспособность при аппаратных сбоях или пиковых нагрузках;
			\item Учет и аудит - возможность детального логирования событий безопасности, формирования отчетов и анализа попыток несанкционированного доступа;
			\item Форм-фактор - физическое исполнение устройства;
			\item Уровень сертификации и класс - наличие сертификатов регуляторов, подтверждающих соответствие требованиям защиты информации для конкретных типов объектов.
	\end{itemize}
	
	\newpage
	
	\section*{Списки доступа и их настройка}
	
	\subsection*{Определение списков контроля доступа}
	ACL (Access Control List) — это набор текстовых выражений, которые что-то разрешают, либо что-то запрещают. Обычно ACL разрешает или запрещает IP-пакеты, но помимо всего прочего он может заглядывать внутрь IP-пакета, просматривать тип пакета, TCP и UDP порты.
	
	\begin{itemize}
		\item ACL могут размещаться на входящем или исходящем направлении.
		\item Обработка ведется по правилам списка доступа строго в том порядке, в котором записаны выражения. Если первое условие не совпало, идет обработка второго условия, если оно совпало, обработка прекращается, если нет, идет обработка третьего условия и так дальше пока не проверятся все. 
		\item Пакет уничтожается, если никакое из условий не совпадает. 
		\item Нельзя разместить более 1 списка доступа на интерфейс, протокол или направление.
		\item ACL не действует на трафик, сгенерированный самим маршрутизатором.
		\item Для фильтрации адресов используется обратная маска.
	\end{itemize}

	\subsection*{Типы списков контроля доступа}
	
		\subsubsection*{Стандартные ACL}
		Проверяют только IP-адрес источника. Нумеруются от 1 до 99 и от 1300 до 1999, также могут быть именованными. Размещаются ближе к получателю.
		
		\begin{minted}[tabsize=4, breaklines=true]{bash}
#Нумерованный список
Router(config)# access-list access-list-number {deny | permit} {source [source-wildcard] | any | host source} [log]
#Именованный список 
Router(config)# ip access-list standard {access-list-name | access-list-number}
Router(config-std-nacl)# [sequence-number] {permit | deny} {source [source-wildcard] | any | host source} [log]

		\end{minted}
		
		\begin{itemize}
			\item \textbf{permit / deny} - разрешить или запретить прохождение трафика.
			\item \textbf{remark} - текстовый комментарий к правилу.
			\item \textbf{address / host / any} - адрес сети, конкретный хост или все адреса.
			\item \textbf{source-wildcard} - обратная маска подсети.
			\item \textbf{log} - включение логирования совпадений с данным правилом.
		\end{itemize}
		
		\begin{minted}[tabsize=4, breaklines=true]{bash}
#Прикрепить ACL к интерфейсу
Router(config-if)#ip access-group <номер списка или имя ACL> {in | out}
		\end{minted}

		\begin{itemize}
			\item \textbf{in} - входящее направление.
			\item \textbf{out} - исходящее направление.

		\end{itemize}
		
		\subsubsection*{Расширенные ACL}
		Фильтрует по IP источника/назначения, протоколу и портам. Размещается ближе к источнику для экономии ресурсов сети. Нумеруются от 100 до 199
		и от 2000 до 2699, также могут быть именованными.
		
		\begin{minted}[tabsize=4, breaklines=true]{bash}
#Нумерованный список
Router(config)# access-list access-list-number {deny | permit} protocol source source-wildcard destination destination-wildcard [precedence precedence] [tos tos] [log | log-input] [time-range time-range-name] [fragments]
#Именованный список
Router(config)# ip access-list extended {access-list-name | access-list-number}
Router(config-ext-nacl)# [sequence-number] {permit | deny} protocol source [source-wildcard] destination [destination-wildcard] [precedence precedence] [tos tos] [log | log-input] [time-range time-range-name] [fragments]
		\end{minted}
		
		\begin{itemize}
			\item \textbf{deny / permit} - разрешить или запретить прохождение трафика, соответствующего критериям.
			\item \textbf{protocol} - сетевой протокол (ip, tcp, udp, icmp, eigrp, ospf и др.) или его номер от 0 до 255.
			\item \textbf{source / source-wildcard} - IP-адрес источника и его обратная маска подсети (также можно использовать ключевые слова \textbf{host} или \textbf{any}).
			\item \textbf{destination / destination-wildcard} - IP-адрес назначения и его обратная маска подсети (также можно использовать ключевые слова \textbf{host} или \textbf{any}).
			\item \textbf{precedence} - фильтрация трафика по значению приоритета.
			\item \textbf{tos} - фильтрация трафика по типу обслуживания.
			\item \textbf{log} - включение логирования совпадений с данным правилом.
			\item \textbf{log-input} - расширенное логирование, которое дополнительно включает в лог входной интерфейс маршрутизатора и MAC-адрес источника.
			\item \textbf{time-range} - привязка правила к созданному ранее временному интервалу.
			\item \textbf{fragments} - применение правила исключительно к фрагментированным IP-пакетам.
		\end{itemize}
		
		\begin{minted}[tabsize=4, breaklines=true]{bash}
#Прикрепить ACL к интерфейсу
Router(config-if)#ip access-group <номер списка или имя ACL> {in | out}
		\end{minted}
		
		\begin{itemize}
			\item \textbf{in} - входящее направление.
			\item \textbf{out} - исходящее направление.
			
		\end{itemize}
		
	\subsubsection*{Динамические}
	Динамические списки доступа используются для удаленных подключений к сети компании. Идея динамических списков доступа состоит в том, что человек должен сначала пройти аутентификацию и только в случае успеха будет применен список доступа, разрешающий доступ к ресурсам сети.
	
	\subsubsection*{Рефлексивные}
	Рефлексивные списки доступа полностью блокируют доступ, но формируется еще один специальный список доступа, который может читать параметры пользовательских сессий, сгенерированных из локальной сети, и для них открывать проход, из внешней сети установить соединение никто не сможет, а на сессии, сгенерированные из локальной сети, будут приходить ответы.
	
	\subsubsection*{Временные}
	Работают,как обычные списки доступа, но с ограничением по времени. Можно ввести специальное расписание, которое активирует ту или иную запись списка доступа.
\end{document}
